\documentclass[11pt]{article}
\usepackage[margin=1in]{geometry}
\usepackage{amsmath,amssymb,amsthm,booktabs,enumitem}
\usepackage[hidelinks]{hyperref}
\newtheorem{definition}{Definition}
\newcommand{\E}{\mathbb{E}}
\newcommand{\R}{\mathbb{R}}
\title{Problem Formulation: Run-Time Detection of Hallucinations\\in Decoder-Free Latent World Models}
\author{}
\date{}

\begin{document}
\maketitle

\section{Setting}

\paragraph{Environment.}
We consider a Markov decision process $\mathcal{M}=(\mathcal{S},\mathcal{A},P,r,\rho_0,\gamma)$ with state space
$\mathcal{S}$, continuous action space $\mathcal{A}\subseteq[-1,1]^{m}$, transition kernel $P(s'\mid s,a)$, reward
function $r:\mathcal{S}\times\mathcal{A}\to\R$, initial-state distribution $\rho_0$ and discount $\gamma\in(0,1)$.
The agent observes $o_t=s_t$ (low-dimensional state observations). Episodes have a fixed length $T$ and do not
terminate early. In our experiments $\mathcal{M}$ is a ManiSkill3 manipulation task (PushCube, StackCube,
PickSingleYCB, PegInsertionSide) with normalized dense reward, simulated deterministically given the full simulator
state.

\paragraph{World model and agent.}
The agent is a trained TD-MPC2 model with parameters $\theta$, consisting of
\begin{align*}
&\text{encoder} && z_t = h_\theta(o_t), &
&\text{latent dynamics} && \hat z' = d_\theta(z,a),\\
&\text{reward head} && \hat r = R_\theta(z,a), &
&\text{critic ensemble} && \{Q_\theta^{(i)}(z,a)\}_{i=1}^{K},\ K=5,\\
&\text{policy prior} && a=\pi_\theta(z), &
&\text{value} && V_\theta(z)=\tfrac1K\textstyle\sum_i Q^{(i)}_\theta\big(z,\pi_\theta(z)\big).
\end{align*}
There is \emph{no decoder}: latents are never trained to reconstruct observations, only to predict rewards and values
and to be self-predictive. $d_\theta$ is deterministic. The model is trained with the standard TD-MPC2 objective
on data the agent collects online; \emph{the detection problem below treats $\theta$ as fixed}.

\paragraph{Imagined rollout.}
Given a real observation $o_t$ and an action sequence $\mathbf a=(a_0,\dots,a_{H-1})$, the model imagines
\[
\hat z_0=h_\theta(o_t),\qquad \hat z_{k+1}=d_\theta(\hat z_k,a_k),\qquad \hat r_k=R_\theta(\hat z_k,a_k),\qquad k=0,\dots,H-1.
\]
The planner (MPPI) scores candidate sequences with imagined rewards plus a terminal value,
$\sum_{k<H}\gamma^k\hat r_k+\gamma^H V_\theta(\hat z_H)$; TD-MPC2 uses $H=3$ and replans at every step.

\section{What counts as a hallucination}

A hallucination is defined only through quantities the world model is trained to predict: rewards along the given
actions. No reconstruction or state probe is involved.

\begin{definition}[Imagination error]
Let $(r_{t},r_{t+1},\dots)$ be the rewards the environment returns when the \emph{same} actions $\mathbf a$ are
executed from the \emph{same} state $s_t$. The imagination error after $k+1$ imagined transitions is
\[
E_k \;=\; \Bigl|\,\sum_{j=1}^{k+1}\gamma^{\,j-1}\bigl(\hat r_j - r_{t+j}\bigr)\Bigr|,\qquad k=0,\dots,H-2,
\]
the absolute difference between the imagined and the real discounted return accumulated from $\hat z_1$ to
$\hat z_{k+1}$. ($\hat r_0$ is excluded because $\hat z_0$ is the encoding of a real observation.)
\end{definition}

\begin{definition}[Hallucination label]
Let $\mathcal{C}$ be a set of calibration windows (Section~3). With $q_\alpha$ the $\alpha$-quantile of the
one-transition errors $\{E_0\}$ over $\mathcal{C}$, set $\varepsilon_{\text{hall}}=q_{0.99}$ and
$\varepsilon_{\text{clean}}=q_{0.90}$. Step $k$ of a window is labelled
\[
y_k=\begin{cases}1 & E_k>\varepsilon_{\text{hall}} \quad\text{(hallucinated)}\\
0 & E_k\le\varepsilon_{\text{clean}} \quad\text{(correct)}\\
\text{excluded} & \text{otherwise.}\end{cases}
\]
A positive step is \emph{sudden} if some increment $E_j-E_{j-1}$, $j\le k$, exceeds the 0.99-quantile of calibration
increments, and \emph{gradual} otherwise.
\end{definition}

The thresholds are relative to the model's own one-step accuracy, so normal approximation and environment noise
are not counted as hallucination. Because $\mathcal{M}$ is simulated deterministically, a single real execution
gives the real return; in a stochastic environment $r_{t+j}$ would be replaced by its expectation over repeated
executions from $s_t$.

\section{The detection task}

\paragraph{Goal.}
Learn a detector
\[
f:\ \bigl(o_t,\ \mathbf a\bigr)\ \longmapsto\ (p_0,\dots,p_{H-2})\in[0,1]^{H-1},
\]
where $p_k$ estimates $\Pr(y_k=1)$: whether the imagined return up to $\hat z_{k+1}$ is hallucinated.

\paragraph{Information available.}
\begin{center}
\begin{tabular}{@{}lll@{}}
\toprule
Phase & Available & Not available\\
\midrule
World-model training & environment interaction, rewards (standard TD-MPC2) & --- \\
Detector calibration & world model $\theta$; $N_{\text{cal}}$ logged real episodes & simulator resets\\
 & $\{(o_{0:T},a_{0:T-1},r_{0:T-1})\}$ collected by the agent & \\
Inference (run time) & world model $\theta$; current $o_t$; candidate actions $\mathbf a$ & real future $o_{t+j}$, $r_{t+j}$; simulator\\
\bottomrule
\end{tabular}
\end{center}
Calibration needs only ordinary logged episodes: for every window of length $H$ inside a logged episode,
the actions and the real rewards are known, so $E_k$ and $y_k$ can be computed without resetting the environment.
At inference the detector may use only the world model's own outputs along $\mathbf a$.

\paragraph{Our detector (for reference).}
With $\bar Q=\frac1K\sum_i Q^{(i)}_\theta$, from the imagined rollout we compute Bellman-consistency residuals, e.g.\ the one-step residual
$\delta_k=\bigl|\bar Q(\hat z_k,a_k)-\bigl(\hat r_k+\gamma V_\theta(\hat z_{k+1})\bigr)\bigr|$ and the residual
anchored at the real start
$\bigl|\bar Q(\hat z_0,a_0)-\bigl(\sum_{j\le k}\gamma^j\hat r_j+\gamma^{k+1}V_\theta(\hat z_{k+1})\bigr)\bigr|$,
plus multi-step, cumulative, bias-corrected and per-head variants, critic disagreement, $V_\theta(\hat z_{k+1})$,
$\hat r_k$ and $k$. A gradient-boosted tree classifier maps these features to $p_k$; it is fitted on calibration
labels $y_k$. No extra models are trained.

\section{Evaluation}

\paragraph{Protocol.}
For each trained world model: collect $N=50$ fresh episodes with the agent's own planner; use the first 25 for
calibration (thresholds and detector fitting) and the other 25 for testing; use all windows of length $H=12$
starting at every real time step, with $\mathbf a$ the actions the agent actually took. Planner-chosen action
sequences are evaluated separately by executing them from saved simulator states.

\paragraph{Metric.}
Detection quality is the step-stratified AUROC on the test windows:
\[
\text{AUROC}_{\text{strat}}=\frac{\sum_k n^+_k n^-_k\,\text{AUROC}_k}{\sum_k n^+_k n^-_k},
\]
where $\text{AUROC}_k$ is computed over steps with label $y_k\in\{0,1\}$ at imagined step $k$ only, and $n^\pm_k$
are the numbers of positives and negatives. Stratifying by $k$ prevents a detector from scoring well merely by
knowing that errors grow with the horizon. It is reported for all, sudden and gradual hallucinations, with 95\%
confidence intervals from a bootstrap over test episodes, and compared with published detectors computed on the
same windows (dynamics-ensemble disagreement, MOBILE-style Bellman-target spread, critic disagreement,
ELVIS-style UCB).

\paragraph{A correct prediction.}
$p_k$ is correct to the extent that it ranks hallucinated steps ($y_k=1$) above correct steps ($y_k=0$) at the same
imagined step $k$. For an operating point, a warning is raised when $p_k$ exceeds a threshold chosen on calibration
data for a target warning rate.

\section{Downstream question (secondary)}
Whether the detector improves decisions is evaluated separately by the task success rate of the agent (fraction of
episodes in which the task is solved at least once), on fixed evaluation seeds shared by all compared methods,
with differences paired by seed.

\end{document}
